% ============================================================================
%  SECCIÓN 6.4: OPERACIONES CRUD CON SUPABASE
% ============================================================================

\section{Operaciones CRUD con Supabase}

En este avance la aplicaci\'on deja de mostrar datos simulados y realiza
operaciones de lectura y escritura directamente sobre la base de datos
centralizada. El resultado es un CRUD completo sobre las entidades
principales del proyecto, conectado a datos reales.

\subsection{Datos de prueba}

Antes de conectar la aplicaci\'on se insertaron registros de prueba para
comprobar que las tablas funcionan correctamente. El \textit{seed}
incluye tres usuarios con distintos roles y un cat\'alogo de escenarios
deportivos reales:

\begin{table}[H]
  \centering
  \caption{Usuarios de prueba del \textit{seed}}
  \contenidoConFuente{%
    \begin{tablaAPA}
      \begin{tabular}{@{}p{4.4cm}p{4.4cm}p{3.4cm}@{}}
        \toprule
        \textbf{Email} & \textbf{Contrase\~na} & \textbf{Rol} \\
        \midrule
        \texttt{admin@test.com} & password123 & admin \\
        \addlinespace
        \texttt{gestor@test.com} & password123 & gestor \\
        \addlinespace
        \texttt{asistente@test.com} & password123 & asistente \\
        \bottomrule
      \end{tabular}
    \end{tablaAPA}%
  }{Elaboraci\'on propia en base al archivo \texttt{seed.sql}.}
\end{table}

\begin{table}[H]
  \centering
  \caption{Ejemplo de registros de escenarios}
  \contenidoConFuente{%
    \begin{tablaAPA}
      \begin{tabular}{@{}p{6.4cm}p{3.0cm}p{3.4cm}@{}}
        \toprule
        \textbf{Nombre} & \textbf{Tipo} & \textbf{Capacidad} \\
        \midrule
        Estadio Hernando Siles & Estadio & 40000 \\
        \addlinespace
        Coliseo Eduardo Leon & Coliseo & 5000 \\
        \addlinespace
        Complejo Deportivo Capriles & Complejo & 3000 \\
        \bottomrule
      \end{tabular}
    \end{tablaAPA}%
  }{Elaboraci\'on propia en base al archivo \texttt{seed.sql}.}
\end{table}

Los registros se verificaron desde el panel de Supabase (Table Editor)
antes de continuar con la conexi\'on desde la aplicaci\'on.

\subsection{El concepto de CRUD}

Las operaciones sobre la base de datos se agrupan en cuatro acciones
b\'asicas conocidas como \textbf{CRUD}:

\begin{enumerate}
  \item \textbf{C - Create (crear):} insertar un nuevo registro.
  \item \textbf{R - Read (leer):} consultar los registros almacenados.
  \item \textbf{U - Update (actualizar):} modificar un registro.
  \item \textbf{D - Delete (eliminar):} eliminar un registro.
\end{enumerate}

El CRUD es uno de los conceptos centrales del desarrollo: todas las
pantallas de gesti\'on de la aplicaci\'on se basan en estas cuatro
operaciones.

\subsection{Flujo de datos din\'amicos}

La pantalla principal ya no muestra datos escritos en el c\'odigo. El
flujo es: la aplicaci\'on solicita los escenarios a Supabase, Supabase
consulta PostgreSQL y los registros regresan para mostrarse en una
\textit{ListView}:

\begin{figure}[H]
\centering
  \figuraAPA{fig:flujo-datos-supabase}{Flujo de consulta de datos din\'amicos}{%
  \begin{tikzpicture}[
    node distance=0.6cm,
    caja/.style={rectangle, draw=black!60, rounded corners,
      minimum width=4.6cm, minimum height=0.65cm, align=center,
      font=\small},
    flecha/.style={-Stealth, thick}
  ]
    \node[caja] (app) {Aplicaci\'on (ListView)};
    \node[caja, below=of app] (sup) {Supabase (consulta)};
    \node[caja, below=of sup] (pg) {PostgreSQL (tablas)};
    \node[caja, right=1.4cm of app] (ui)
      {Escenarios\\ \footnotesize desde la base de datos};
    \draw[flecha] (app) -- node[right, font=\footnotesize]
      {solicita} (sup);
    \draw[flecha] (sup) -- node[right, font=\footnotesize]
      {SELECT} (pg);
    \draw[flecha] (pg) -- node[left, font=\footnotesize]
      {registros} (sup);
    \draw[flecha] (sup) -- node[left, font=\footnotesize]
      {devuelve} (app);
    \draw[flecha] (app) -- node[above, font=\footnotesize]
      {muestra} (ui);
  \end{tikzpicture}
  }{Elaboraci\'on propia.}
\end{figure}

\subsection{Implementaci\'on del CRUD}

Las operaciones se organizan en \textit{hooks} de React Query que
centralizan las consultas y el manejo de estados. A continuaci\'on se
describen las implementaciones de cada operaci\'on del CRUD.

\subsubsection{Read: consulta de escenarios}

El hook \texttt{useScenarios} lee los escenarios activos con sus im\'agenes
y deportes, ordenados por nombre:

\begin{lstlisting}
async function fetchScenarios() {
  const { data, error } = await supabase
    .from('scenarios')
    .select(`
      *,
      scenario_images(url, is_primary, storage_path, display_order),
      scenario_sports(sports(nombre))
    `)
    .eq('estado', 'activo')
    .order('nombre');

  if (error) throw new Error(error.message);
  return data;
}

export function useScenarios() {
  return useQuery({
    queryKey: ['scenarios'],
    queryFn: fetchScenarios,
    staleTime: 1000 * 60 * 5,
  });
}
\end{lstlisting}

El resultado alimenta las pantallas de cat\'alogo y del mapa, por lo que
los datos mostrados provienen directamente de la base de datos.

\subsubsection{Create: creaci\'on de registros}

La creaci\'on de registros se realiza con \texttt{.insert()}. Por ejemplo,
guardar un favorito:

\begin{lstlisting}
async function addFavorite(userId: string, scenarioId: string) {
  const { error } = await supabase
    .from('favorites')
    .insert({ user_id: userId, scenario_id: scenarioId });

  if (error) throw new Error(error.message);
}
\end{lstlisting}

En el caso de los escenarios, el formulario de gesti\'on valida los
campos con Zod y ejecuta la inserci\'on; al finalizar se invalida el
cach\'e de React Query para que la lista se actualice.

\subsubsection{Update: actualizaci\'on de registros}

La actualizaci\'on se realiza con \texttt{.update()} sobre el registro
identificado. El formulario de edici\'on reutiliza el mismo componente
del alta, cargando los datos existentes y enviando los nuevos valores:

\begin{lstlisting}
const { error } = await supabase
  .from('scenarios')
  .update({ capacidad: nuevaCapacidad, descripcion: nuevaDescripcion })
  .eq('id', escenarioId);
\end{lstlisting}

\subsubsection{Delete: eliminaci\'on de registros}

La eliminaci\'on se implementa de dos formas seg\'un la entidad:

\begin{itemize}
  \item \textbf{Eliminaci\'on suave (soft delete):} los escenarios no se
    borran de la base de datos; su estado cambia a ``inactivo'' y
    desaparecen del cat\'alogo p\'ublico.
  \item \textbf{Eliminaci\'on f\'isica:} eventos y favoritos se eliminan
    directamente con \texttt{.delete()}.
\end{itemize}

\begin{lstlisting}
async function removeFavorite(userId: string, scenarioId: string) {
  const { error } = await supabase
    .from('favorites')
    .delete()
    .eq('user_id', userId)
    .eq('scenario_id', scenarioId);

  if (error) throw new Error(error.message);
}
\end{lstlisting}

La subida de im\'agenes a Supabase Storage, documentada en el avance 5,
complementa el CRUD: cada imagen se sube al bucket \texttt{scenario-images}
y su referencia se registra en la tabla \texttt{scenario\_images}.

\subsection{Conexi\'on del CRUD con el proyecto real}

El CRUD no es un m\'odulo gen\'erico separado del proyecto: cada
operaci\'on corresponde a una funcionalidad de \nombreApp{} definida desde
el inicio. La relaci\'on es la siguiente:

\begin{itemize}
  \item \textbf{CRUD de escenarios:} creaci\'on, consulta, edici\'on y
    baja de escenarios deportivos (gesti\'on interna para admin y gestor).
  \item \textbf{CRUD de eventos:} creaci\'on y eliminaci\'on de eventos en
    cada escenario.
  \item \textbf{CRUD de favoritos:} consultar, agregar y quitar
    escenarios favoritos del usuario autenticado.
\end{itemize}

\subsection{Resumen de operaciones CRUD}

\begin{table}[H]
  \centering
  \caption{Resumen de operaciones CRUD implementadas}
  \contenidoConFuente{%
    \begin{tablaAPA}
      \begin{tabular}{@{}p{2.6cm}p{3.6cm}p{5.4cm}@{}}
        \toprule
        \textbf{Entidad} & \textbf{Operaci\'on} & \textbf{Hook / funci\'on} \\
        \midrule
        escenarios & Read & {\footnotesize\ttfamily useScenarios /
          useScenario} \\
        \addlinespace
        escenarios & Create / Update & {\footnotesize\ttfamily
          formulario (upsert)} \\
        \addlinespace
        escenarios & Delete & {\footnotesize\ttfamily soft delete
          (estado = inactivo)} \\
        \addlinespace
        eventos & Create / Delete & {\footnotesize\ttfamily
          rutas de evento del gestor} \\
        \addlinespace
        favoritos & Create / Read / Delete & {\footnotesize\ttfamily
          useToggleFavorite} \\
        \bottomrule
      \end{tabular}
    \end{tablaAPA}%
  }{Elaboraci\'on propia en base al c\'odigo fuente del proyecto.}
\end{table}

Con estas operaciones la aplicaci\'on queda conectada a datos reales:
lo que se crea desde un dispositivo queda disponible para el resto de los
usuarios a trav\'es de la base de datos centralizada.